\section{Loi générale des masses : atlas généalogique, lecteurs bidirectionnels et droite dimensionnelle}
\label{10j-sec-ley-general-masas}

La loi générale des masses n'est pas une table de valeurs ni une règle d'ajustement par
espèce. C'est la réalisation dimensionnelle d'une généalogie construite avant
d'ouvrir tout catalogue métrologique :
\begin{equation}
 \APP\longrightarrow\TRIT\longrightarrow\TPK
 \longrightarrow\widetilde\Gamma^{\rm surv}
 \longrightarrow\Gamma^{\rm obs}
 \longrightarrow L_{108}
 \longrightarrow\sigma\in\mathbb Z^5
 \longrightarrow\chi_0
 \longrightarrow\chi_{\rm full}
 \longrightarrow\widehat{\mathcal M}
 \longrightarrow\mathcal L_M .
 \label{10j-eq-cadena-masas}
\end{equation}
APP apporte les deux feuillets arithmétiques et leur défaut orienté ; TRIT conserve
signe, seuil et régime ; le TPK sélectionne, transporte et met à jour l'état
au moyen de
\(
 U_t=\operatorname{Upd}_t\circ\operatorname{Tra}_t\circ
 \operatorname{Sel}_t
\), en conservant résidu, quotient, retenue, feuillet, orientation, bord,
voie, mémoire et survie. La grandeur physique n'apparaît qu'à la fin
de \eqref{10j-eq-cadena-masas}. En particulier, ni une masse connue ni une
étiquette d'espèce n'interviennent dans la sélection de la voie.

Les constantes qui apparaissent ci-dessous sont reçues comme publications internes de
la même généalogie HMT. En particulier, \(\pi_{\rm HMT}\), \(e_{\rm HMT}\),
\(\varphi_{\rm HMT}\) et \(\alpha_{\rm HMT}\) appartiennent à l'image nonadique
typée engendrée par APP--TRIT--TPK. Leur utilisation ultérieure dans la branche de masse est une
composition interne descendante ; ce n'est pas une injection de valeurs
conventionnelles.

\subsection{Domaine fini, voies et multisections}

\begin{definition}[Atlas orienté et quotients de fibre]
\label{10j-def-atlas}
Soit \(\mathcal C_{81}\) l'atlas cellulaire produit par le lecteur de familles
de l'état enrichi. Nous définissons
\begin{equation}
 \mathcal R_{324}:=
 \mathcal C_{81}\times\{+1,-1\}\times\{+1,-1\}.
 \label{10j-eq-rutas-324}
\end{equation}
Le premier quotient identifie les cellules ayant la même histoire non orientée et donne
l'ensemble \(\mathcal F_{56}\) des familles de voie. La stratification par
type de représentation donne treize multisections
\(\mathcal M_{13}\). La carte nulle \(\mathcal N\) est un terme séparé du
domaine positif ; ce n'est pas une quatorzième multisection obtenue en annulant
un caractère.
\end{definition}

\begin{theorem}[Recensement complet et absence de sélection ponctuelle]
\label{10j-thm-dominio}
Le domaine interne de la loi satisfait
\begin{equation}
 |\mathcal C_{81}|=81,
 \qquad |\mathcal F_{56}|=56,
 \qquad |\mathcal M_{13}|=13,
 \qquad |\mathcal R_{324}|=324.
 \label{10j-eq-censo-dominio}
\end{equation}
Les multiplicités du quotient des familles sont
\begin{equation}
 41\times1,\quad10\times2,\quad2\times3,\quad1\times4,\quad2\times5,
 \qquad 41+20+6+4+10=81,
 \label{10j-eq-familias-56}
\end{equation}
et les rangs des treize multisections sont
\begin{equation}
 \underbrace{1}_{1\ \text{ fois}},\quad
 \underbrace{2}_{2\ \text{ fois}},\quad
 \underbrace{4}_{3\ \text{ fois}},\quad
 \underbrace{6}_{2\ \text{ fois}},\quad
 \underbrace{8}_{3\ \text{ fois}},\quad
 \underbrace{12}_{1\ \text{ fois}},\quad
 \underbrace{16}_{1\ \text{ fois}},
 \label{10j-eq-rangos-13}
\end{equation}
dont la somme avec multiplicité est \(81\). La fibre de rang un est la fibre
électronique centrale. Toute fibre non centrale est de rang supérieur à un et, par
conséquent, sa sortie naturelle est un opérateur sur la fibre ; un nom physique ne
sélectionne pas à lui seul l'une de ses valeurs propres.
\end{theorem}

\begin{proof}
Chaque cellule admet exactement les quatre orientations indépendantes
\((h_0,v_0)\in\{\pm1\}^2\), d'où
\(|\mathcal R_{324}|=4\cdot81=324\). Les sommes de
\eqref{10j-eq-familias-56} donnent simultanément cinquante-six classes et
quatre-vingt-une cellules. La somme des rangs de \eqref{10j-eq-rangos-13} est
\[
 1+2\cdot2+3\cdot4+2\cdot6+3\cdot8+12+16=81.
\]
L'inversion cellulaire possède un unique point fixe, le centre \((5,5)\) ; toutes les autres orbites forment des fibres de cardinal supérieur à un. Un choix ponctuel invariant dans l'une de ces orbites contredirait la transitivité de l'action d'inversion sur la fibre. C'est pourquoi seul le centre admet canoniquement une réduction de rang un.
\end{proof}

\subsection{Du défaut APP au registre entier}

\begin{proposition}[Séparation typée de \(4\) et \(54\)]
\label{10j-prop-cuatro-cincuentaycuatro}
Soit \(R\) l'involution qui échange les feuillets APP et \(P_\pm=(I\pm R)/2\). Le bloc central
\begin{equation}
 B_c=7I+2R=9P_++5P_-
 \label{10j-eq-bloque-app}
\end{equation}
produit le saut spectral primordial \(9-5=4\). Les nombres
\begin{equation}
 2,\qquad 6=51-45,\qquad
 54=9\cdot6=459-405
 \label{10j-eq-despliegue-2-6-54}
\end{equation}
sont respectivement le couplage local, le défaut de compression et le défaut intégré. Dans le lecteur électronique, \(54=D_1(\Gamma_e)\) compte en outre un retour cinématique. Aucune de ces occurrences ne remplace le saut APP \(4\).
\end{proposition}

\begin{proof}
La décomposition spectrale de \eqref{10j-eq-bloque-app} donne les valeurs propres \(9\) et \(5\), donc leur différence est \(4\). L'entrée hors diagonale dans la base des feuillets est \(2\). Les deux compressions globales s'évaluent en \(51\) et \(45\), et leur différence est \(6\) ; intégrer ce défaut sur neuf positions donne \(54\). Le domaine, l'opération et le codomaine changent à chaque étape. L'égalité numérique ultérieure avec \(D_1(\Gamma_e)\) est une confluence de lecteurs dans \(\mathbb Z\), non une identité des objets qu'ils évaluent.
\end{proof}

Pour une voie \(\Gamma\in\mathcal R_{324}\), le registre CT--108 conserve un mot temporel \(w_\Gamma\in\mathbb F_3^{108}\), un mot signé \(\tau_\Gamma\in\{-1,0,+1\}^{12}\), le feuillet, l'enroulement, la mémoire et la retenue. Le retour nonadique satisfait
\begin{equation}
 g(K+9)=g(K),\qquad B_{K+9}\ne B_K\quad\hbox{en général},
 \label{10j-eq-retorno-memoria}
\end{equation}
car la phase revient et le registre dodécaphasique avance.

\begin{definition}[Signature entière et caractère central]
\label{10j-def-firma-caracter}
Le registre de la voie définit
\begin{equation}
 \sigma(\Gamma)=
 (n_A,s_C,k_{\Delta_4},b_{120},b_\zeta)\in\mathbb Z^5.
 \label{10j-eq-firma-z5}
\end{equation}
La fermeture de douze événements détermine la lecture orientée \(s_C^{(12)}=s_C/6\) ; le diviseur six n'est pas un paramètre ajusté. Soient
\begin{equation}
 \Delta_4:=
 \frac{\pi_{\rm HMT}-e_{\rm HMT}\log\pi_{\rm HMT}}{270},
 \qquad \zeta_{270}:=135\Delta_4^2,
 \label{10j-eq-delta-zeta}
\end{equation}
\begin{equation}
 A:=1000\alpha_{\rm HMT},qquad
 x:=\frac{\pi A}{180},qquad
 y:=\frac{\pi C^*}{180},qquad
 q_\pm:=e^{-x\mp y},qquad
 s_n:=q_-^n-q_+^n .
 \label{10j-eq-canales}
\end{equation}
Ici \(C^*\) est le contre-angle déjà construit par le sélecteur régional HMT, non par une masse. Le caractère central est
\begin{equation}
 \chi_0(\sigma):=
 \exp\!\left(
 n_Ax+\frac{s_C}{6}y+k_{\Delta_4}\Delta_4
 +b_{120}s_{120}+b_\zeta\zeta_{270}
 \right).
 \label{10j-eq-caracter-central}
\end{equation}
\end{definition}

\begin{lemma}[Homomorphisme du caractère]
\label{10j-lem-caracter}
Pour des signatures composables, \(
 \chi_0(\sigma_1+\sigma_2)=\chi_0(\sigma_1)\chi_0(\sigma_2)
\). Le caractère est positif, adimensionnel et ne contient aucune unité de masse.
\end{lemma}

\begin{proof}
L'exposant de \eqref{10j-eq-caracter-central} est une forme linéaire sur \(\mathbb Z^5\) ; l'exponentielle transforme la somme en produit. Tous ses coefficients sont des coordonnées adimensionnelles construites avant l'ouverture d'une carte dimensionnelle. L'unité n'apparaît qu'au passage à \(\mathcal L_M\).
\end{proof}

\subsection{Échelle absolue d'action et correcteur relatif}

\begin{theorem}[Origine déterminantale de la décennie \(-34\)]
\label{10j-thm-escala-menos34}
Pour
\begin{equation}
 H_4=
 \begin{pmatrix}
 1&1&1&1\\
 1&1&-1&-1\\
 1&-1&1&-1\\
 1&-1&-1&1
 \end{pmatrix}
 \label{10j-eq-h4}
\end{equation}
on a
\begin{equation}
 \operatorname{SNF}(H_4)=\operatorname{diag}(1,2,2,4),
 \qquad |\operatorname{coker}H_4|=16.
 \label{10j-eq-snf-h4}
\end{equation}
Le lecteur déterminantal publie
\begin{equation}
 \Sigma_H:=\varphi_{\rm HMT}\alpha_{\rm HMT}^{16}
 =1.0462438381047565801842292083\ldots\times10^{-34},
 \label{10j-eq-sigma-h}
\end{equation}
d’où \(\operatorname{ord}_{10}(\Sigma_H)=-34\).
\end{theorem}

\begin{proof}
Les diviseurs déterminantaux de \(H_4\) sont \(\delta_1=1\), \(\delta_2=2\), \(\delta_3=4\) et \(\delta_4=|\det H_4|=16\). Par le théorème de Smith, les facteurs invariants sont \(
 d_1=1,
 d_2=\delta_2/\delta_1=2,
 d_3=\delta_3/\delta_2=2,
 d_4=\delta_4/\delta_3=4
\), ce qui prouve \eqref{10j-eq-snf-h4}. La substitution des sorties HMT déjà générées \(\varphi_{\rm HMT}\) et \(\alpha_{\rm HMT}\) dans le lecteur donne la valeur de \eqref{10j-eq-sigma-h} ; en particulier \(10^{-34}<\Sigma_H<10^{-33}\). Aucune constante SI n'a été utilisée pour sélectionner l'exposant seize : celui-ci est l'ordre du conoyau.
\end{proof}

Les sections antérieure et postérieure au retour vivent dans une unique droite d'action :
\begin{equation}
 \hbar_{\rm pre}=H_5\,10^{-34}\mathcal U_S,
 \qquad
 \hbar_{\rm ret}=\eta_{\rm ret}\,10^{-34}\mathcal U_S,
 \label{10j-eq-secciones-accion}
\end{equation}
avec
\begin{align}
 H_5&=1.0545718176461544696258218294330594765\ldots,\notag\\
 \eta_{\rm ret}&=1.0545718169150390611336905847242997339\ldots,\notag\\
 R_{\rm act}&:=\frac{H_5}{\eta_{\rm ret}}
 =1.000000000693281763048512336828598947\ldots .
 \label{10j-eq-ract}
\end{align}
La différence \(H_5-\eta_{\rm ret}\) et le quotient \(R_{\rm act}\) sont les lectures additive et multiplicative des deux mêmes sections, non deux constantes de Planck.

\begin{proposition}[Séparation de l'échelle et du raffinement]
\label{10j-prop-escala-corrector}
Soit \(t_0>0\) la section temporelle déclarée et
\begin{equation}
 \omega_0=\frac{2\pi_{\rm HMT}}{108t_0},
 \qquad \varepsilon_{\rm clk}=\hbar_{\rm ret}\omega_0,
 \qquad m_{\rm clk}=\varepsilon_{\rm clk}c_{\rm int}^{-2}.
 \label{10j-eq-reloj-masa}
\end{equation}
Alors \(10^{-34}\mathcal U_S\) fixe la section absolue de la droite de masse, tandis que \(R_{\rm act}^{\kappa}\) est un correcteur relatif positif. La décennie \(-34\) ne fait pas partie de \(\kappa\) et s'annule exactement dans tout rapport de masses.
\end{proposition}

\begin{proof}
Par \eqref{10j-eq-secciones-accion}, les deux sections contiennent le même facteur \(10^{-34}\mathcal U_S\). Dans le quotient, on obtient \(\hbar_{\rm pre}/\hbar_{\rm ret}=H_5/\eta_{\rm ret}=R_{\rm act}\), si bien que la puissance décimale s'annule. Si \(m_X=m_{\rm clk}\chi_XR_{\rm act}^{\kappa_X}\) et \(m_Y=m_{\rm clk}\chi_YR_{\rm act}^{\kappa_Y}\), alors
\[
 \frac{m_Y}{m_X}
 =\frac{\chi_Y}{\chi_X}
  R_{\rm act}^{\kappa_Y-\kappa_X}.
\]
Introduire \(-34\) dans \(\kappa\), ou multiplier de nouveau une coordonnée déjà exprimée en MeV par \(10^{-34}\), compterait deux fois la même échelle et violerait le typage.
\end{proof}

\subsection{Lecteurs bidirectionnels et recensement exact}

Pour \(k\in\mathbb Z/108\mathbb Z\), définissons la discordance cyclique
\begin{equation}
 D_\Gamma(k):=
 \sum_{t=0}^{107}
 \mathbf1_{\{w_{\Gamma,t}\ne w_{\Gamma,t+k}\}}.
 \label{10j-eq-discrepancia}
\end{equation}
Le lecteur chronologique et le lecteur dodécaphasique sont
\begin{align}
 K_D(\Gamma)
 &:=\left(D_{27}+D_{36},D_1,\frac{D_4-D_3}{2}\right),
 \label{10j-eq-kd}\\
 K_\Omega(\Gamma)
 &:=(\Omega_{90\mid120}(\tau_\Gamma),54,P_6(\tau_\Gamma)).
 \label{10j-eq-komega}
\end{align}
Le premier lit les déplacements de \(w_\Gamma\) ; le second, l'autocorrélation des fenêtres de \(\tau_\Gamma\). Ils partagent la voie, mais non le quotient d'information.

\begin{theorem}[Profils certifiés des lecteurs]
\label{10j-thm-censo-lectores}
Sur \(\mathcal R_{324}\), \(K_D\) possède exactement quatorze profils, \(K_\Omega\) neuf, et la paire \((K_D,K_\Omega)\) quarante. Il existe dix-huit voies sur lesquelles les deux triplets coïncident et l'unique triplet coïncident est
\begin{equation}
 K_D(\Gamma)=K_\Omega(\Gamma)=(80,54,6).
 \label{10j-eq-unica-coincidencia}
\end{equation}
Les multiplicités marginales sont celles des \cref{10j-tab-kd,10j-tab-komega}.
\end{theorem}

\begin{longtable}{@{}rrr r@{}}
\caption{Quatorze profils de \(K_D\) sur les 324 voies.}
\label{10j-tab-kd}\\
\toprule
\(A_D\)&\(B_D\)&\(C_D\)&multiplicité\\
\midrule
\endfirsthead
\toprule
\(A_D\)&\(B_D\)&\(C_D\)&multiplicité\\
\midrule
\endhead
0&24&0&18\\
40&24&2&36\\
40&54&6&18\\
40&78&27&36\\
40&84&30&36\\
60&78&25&18\\
60&84&28&18\\
80&54&6&18\\
80&54&7&18\\
80&66&2&18\\
80&66&3&36\\
80&66&4&18\\
100&66&0&18\\
100&66&2&18\\
\midrule
\multicolumn{3}{r}{total}&324\\
\bottomrule
\end{longtable}

\begin{table}[htbp]
\centering
\caption{Neuf profils de \(K_\Omega\) sur les 324 voies.}
\label{10j-tab-komega}
\begin{tabular}{@{}rrr r@{}}
\toprule
\(\Omega\)&\(B_\Omega\)&\(P_6\)&multiplicité\\
\midrule
80&54&6&198\\
56&54&5&68\\
48&54&5&34\\
36&54&4&8\\
20&54&4&4\\
4&54&2&4\\
40&54&4&4\\
24&54&4&2\\
24&54&2&2\\
\midrule
\multicolumn{3}{r}{total}&324\\
\bottomrule
\end{tabular}
\end{table}

\begin{proof}[Démonstration de \cref{10j-thm-censo-lectores}]
On parcourt le domaine fini de \eqref{10j-eq-rutas-324}. Pour chaque \((c,h_0,v_0)\), l'itération de \(U_t\) pendant 108 étapes détermine de manière unique \(w_\Gamma\) et \(\tau_\Gamma\) ; puis on évalue \eqref{10j-eq-discrepancia}--\eqref{10j-eq-komega}. Le regroupement des triplets égaux produit les quatorze lignes de \cref{10j-tab-kd} et les neuf de \cref{10j-tab-komega} ; les deux sommes valent \(324\). Le regroupement des paires ordonnées donne quarante classes. Enfin, le filtrage des lignes présentant l'égalité des deux triplets laisse dix-huit voies, toutes dans l'unique intersection \((80,54,6)\). L'algorithme ne reçoit que la cellule, l'orientation et la transition TPK ; il ne contient ni colonne de masses ni étiquette PDG.
\end{proof}

Les évaluations scalaires des lecteurs, lorsqu'il convient de les appliquer, sont
\begin{align}
 \kappa_D(\Gamma)
 &=A_D(\Gamma)-\alpha_{\rm HMT}B_D(\Gamma)
   +C_D(\Gamma)\Delta_4,
 \label{10j-eq-kappa-d}\\
 \kappa_\Omega(\Gamma)
 &=\Omega(\Gamma)-54\alpha_{\rm HMT}
   +P_6(\Gamma)\Delta_4.
 \label{10j-eq-kappa-omega}
\end{align}
La présence de \(\alpha_{\rm HMT}\) ici est descendante : la constante a déjà été publiée par HMT avant que le lecteur n'agisse sur une voie.

\subsection{Électron basal et fermeture bêta sans cible}

\begin{theorem}[Fibre électronique et valeur directrice]
\label{10j-thm-electron}
L'unique point fixe central détermine
\begin{equation}
 \Gamma_e=(5,5,++),
 \qquad \tau_e=+++---+++---,
 \qquad
 (D_1,D_3,D_4,D_{27},D_{36})=(54,56,68,48,32).
 \label{10j-eq-ruta-electron}
\end{equation}
Par conséquent,
\begin{equation}
 K_D(\Gamma_e)=K_\Omega(\Gamma_e)=(80,54,6).
 \label{10j-eq-lectores-electron}
\end{equation}
Le sélecteur basal interne \((-22,+15)\) produit
\begin{equation}
\begin{aligned}
 e_0&:=\frac{\sqrt3}{4}
 \left[
  \exp\!\left(\frac{\varphi_{\rm HMT}}{\pi_{\rm HMT}^2}\right)
  -22\alpha_{\rm HMT}^3
 \right](1+15\Delta_4)\\
 &=0.51099892248806607250987087719134943763\ldots\ {\rm MeV}.
\end{aligned}
 \label{10j-eq-electron-basal}
\end{equation}
Avec
\begin{equation}
 \kappa_e:=80-54\alpha_{\rm HMT}+6\Delta_4,
 \label{10j-eq-kappa-electron}
\end{equation}
la sortie bêta est
\begin{equation}
 \boxed{
 m_e^\beta=e_0R_{\rm act}^{\kappa_e}
 =0.5109989506900008086082571648379233956246348677151\ldots\ {\rm MeV}.}
 \label{10j-eq-electron-beta}
\end{equation}
\end{theorem}

\begin{proof}
De \eqref{10j-eq-ruta-electron},
\[
 D_{27}+D_{36}=48+32=80,
 \qquad D_1=54,
 \qquad (D_4-D_3)/2=(68-56)/2=6,
\]
ce qui prouve la première égalité de \eqref{10j-eq-lectores-electron} ; l'autocorrélation du mot dodécaphasique donne le même triplet. Les entiers du sélecteur sont construits sans masse cible :
\begin{equation}
 22=396/18=27-5=24-2,
 \qquad 15=270/18=\binom62=\binom64,
 \label{10j-eq-genealogia-22-15}
\end{equation}
et la matrice locale des lacunes \(
 V_e=\left(\begin{smallmatrix}21&22\\6&7\end{smallmatrix}\right)
\)
a pour déterminant \(15\). La substitution des publications internes dans \eqref{10j-eq-electron-basal} donne la valeur basale indiquée. Le triplet commun \((80,54,6)\) détermine \eqref{10j-eq-kappa-electron} ; enfin, \eqref{10j-eq-ract} donne \eqref{10j-eq-electron-beta}. La valeur externe n'est ouverte qu'après cette évaluation, si bien qu'elle ne participe au choix ni de \((-22,+15)\), ni de la voie, ni de la puissance.
\end{proof}

La valeur basale \(e_0\) est une étape et ne remplace pas \(m_e^\beta\). On ne confond pas non plus la signature brute \((24,0,12,0,-12)\), la signature paire CT--108 \((-12,-4,12,-6,12)\), l'origine affine \(v_e=0\), la voie \(\Gamma_e\) et le mot utilisé par le lecteur d'Euler : ils appartiennent à des types distincts.

\subsection{Semi-groupe global et convergence des tours}

\begin{theorem}[Semi-groupe des hauteurs]
\label{10j-thm-semigrupo}
Les hauteurs de tour forment
\begin{equation}
 \mathfrak S=\langle90,120\rangle
 =\{90,120\}\cup\{30r:r\ge6\}.
 \label{10j-eq-semigrupo}
\end{equation}
Son algèbre de semi-groupe admet la présentation
\begin{equation}
 k[\mathfrak S\cup\{0\}]
 \simeq k[X_{90},Y_{120}]/(X_{90}^4-Y_{120}^3).
 \label{10j-eq-algebra-semigrupo}
\end{equation}
La hauteur \(360\) conserve deux généalogies avant le quotient : quatre étapes de \(90\) et trois de \(120\).
\end{theorem}

\begin{proof}
Diviser par \(30\) réduit le problème au semi-groupe \(\langle3,4\rangle\). Ses éléments positifs sont \(3,4\) et tous les entiers \(r\ge6\) : selon le résidu de \(r\) modulo trois, on écrit \(r=3a+4b\) avec \(b\in\{0,1,2\}\) et \(a\ge0\). L'unique relation primitive entre les deux générateurs est \(4\cdot3=3\cdot4\), qui, lorsque les hauteurs sont rétablies, produit \(X_{90}^4=Y_{120}^3\). Avant d'imposer cette relation, les mots \(X_{90}^4\) et \(Y_{120}^3\) retiennent des histoires distinctes.
\end{proof}

Pour \(h=30r\), nous définissons les coefficients relatifs de tour
\begin{align}
 \eta_{30r}^D(\Gamma)
 &:=D_\Gamma(9r)-D_{\Gamma_e}(9r),
 \label{10j-eq-eta-d}\\
 \eta_{30r}^\Omega(\Gamma)
 &:=A_\Gamma(r)-A_{\Gamma_e}(r),
 \qquad
 A_\Gamma(r):=
 \sum_{j=0}^{11}
 \left(\sum_{u=0}^{r-1}\tau_{\Gamma,j+u}\right)^2,
 \label{10j-eq-eta-omega}
\end{align}
où les indices de \(\tau\) sont pris modulo douze. Le caractère raffiné est
\begin{equation}
 \chi_{\rm full}^{\,r}(\sigma,\eta)
 :=\chi_0(\sigma)
 \exp\!\left(\sum_{h\in\mathfrak S}\eta_h^r s_h\right),
 \qquad r\in\{D,\Omega\}.
 \label{10j-eq-caracter-full}
\end{equation}

\begin{proposition}[Convergence absolue]
\label{10j-prop-convergencia-torres}
Les deux séries de \eqref{10j-eq-caracter-full} convergent absolument pour toute voie. En conséquence, la hiérarchie n'a pas à être tronquée aux huit hauteurs habituellement présentées.
\end{proposition}

\begin{proof}
Par définition, \(0\le D_\Gamma(k)\le108\), donc \(|\eta_h^D|\le108\). Comme chaque somme partielle d'un mot signé est de module au plus \(r\), on obtient \(|\eta_{30r}^\Omega|\le24r^2\). De plus, \(0<q_+<q_-<1\) et \(|s_{30r}|\le2q_-^{30r}\). Par comparaison avec une série géométrique et avec \(\sum r^2q_-^{30r}\), les deux séries convergent absolument.
\end{proof}

La tour est complète en tant que codomaine d'expansion : ce n'est pas un sélecteur physique d'espèce. Un algorithme qui reçoit un rapport pour le développer ne démontre pas quelle voie a produit ce rapport.

\subsection{Les treize paires d'opérateurs de fibre}

Pour chaque multisection \(M\), soit \(
 \mathcal H_M=\bigoplus_{\Gamma\in M}\mathbb C e_\Gamma
\). Nous définissons
\begin{equation}
 B_M^D e_\Gamma:=\kappa_D(\Gamma)e_\Gamma,
 \qquad
 B_M^\Omega e_\Gamma:=\kappa_\Omega(\Gamma)e_\Gamma.
 \label{10j-eq-operadores-bm}
\end{equation}
Les treize fibres canoniques et leurs rangs sont
\begin{equation}
 \begin{array}{c|ccccccccccccc}
 M&\ell_0&\ell_2&\ell_4&\nu_1&\nu_2&\nu_3&\nu_4
  &d_1&d_2&d_3&d_4&u_1&u_3\\ \hline
 \dim M&1&8&16&2&4&6&8&2&4&6&8&4&12.
 \end{array}
 \label{10j-eq-trece-fibras}
\end{equation}

\begin{theorem}[Loi opératorielle sur les multisections]
\label{10j-thm-ley-operatoria}
Les vingt-six opérateurs de \eqref{10j-eq-operadores-bm} sont autoadjoints. Pour un opérateur basal positif \(\widehat M_M^{(0)}\),
\begin{equation}
 \widehat M_M^{\beta,r}
 =R_{\rm act}^{B_M^r/2}
  \widehat M_M^{(0)}
  R_{\rm act}^{B_M^r/2},
 \qquad r\in\{D,\Omega\},
 \label{10j-eq-ley-operatoria}
\end{equation}
est positif et covariant sous changement unitaire de base. Sur la fibre \(\ell_0\) de rang un, il se réduit au scalaire électronique \eqref{10j-eq-electron-beta}. Sur une fibre de rang supérieur, le spectre de \eqref{10j-eq-ley-operatoria} est la sortie complète tant qu'une image directe physique prospective et une scalarisation ne sont pas fixées.
\end{theorem}

\begin{proof}
Dans la base des voies, \(B_M^D\) et \(B_M^\Omega\) sont diagonaux à entrées réelles, donc autoadjoints. Comme \(R_{\rm act}>0\), le calcul fonctionnel définit \(R_{\rm act}^{B_M^r/2}\) comme un opérateur positif et inversible. Pour tout \(v\),
\[
 \langle v,\widehat M_M^{\beta,r}v\rangle
 =\left\langle R_{\rm act}^{B_M^r/2}v,
 \widehat M_M^{(0)}R_{\rm act}^{B_M^r/2}v\right\rangle\ge0.
\]
Si \(U\) est unitaire, le changement \(B\mapsto UBU^*\), \(\widehat M^{(0)}\mapsto
U\widehat M^{(0)}U^*\) porte la sortie à \(U\widehat M^{\beta,r}U^*\) et conserve le spectre, la trace et le déterminant. Lorsque \(\dim\ell_0=1\), les deux lecteurs valent \((80,54,6)\) et la formule est scalaire. Pour \(\dim M>1\), choisir une valeur propre sans morphisme physique supplémentaire n'est pas une opération naturelle sur la fibre.
\end{proof}

Le tableau suivant présente les traces vectorielles normalisées des triplets avant l'application de \(\alpha_{\rm HMT}\) et \(\Delta_4\) ; leur diversité prouve qu'il n'existe pas de paire angulaire unique qui remplace les lecteurs de voie.
\begin{longtable}{@{}c r c c@{}}
\caption{Traces normalisées de \(K_D\) et \(K_\Omega\) sur les treize fibres.}
\label{10j-tab-trazas-13}\\
\toprule
\(M\)&\(\dim M\)&\(\overline{\mathbf K}_M^D\)&
\(\overline{\mathbf K}_M^\Omega\)\\
\midrule
\endfirsthead
\toprule
\(M\)&\(\dim M\)&\(\overline{\mathbf K}_M^D\)&
\(\overline{\mathbf K}_M^\Omega\)\\
\midrule
\endhead
\(\ell_0\)&1&\((80,54,6)\)&\((80,54,6)\)\\
\(\ell_2\)&8&\((65,249/4,17/2)\)&\((80,54,6)\)\\
\(\ell_4\)&16&\((225/4,489/8,21/2)\)&\((141/2,54,11/2)\)\\
\(\nu_1\)&2&\((40,51,29/2)\)&\((60,54,5)\)\\
\(\nu_2\)&4&\((45,63,29/2)\)&\((61,54,5)\)\\
\(\nu_3\)&6&\((170/3,55,34/3)\)&\((206/3,54,11/2)\)\\
\(\nu_4\)&8&\((105/2,237/4,95/8)\)&\((143/2,54,45/8)\)\\
\(d_1\)&2&\((40,51,29/2)\)&\((60,54,5)\)\\
\(d_2\)&4&\((45,63,57/4)\)&\((61,54,5)\)\\
\(d_3\)&6&\((170/3,55,23/2)\)&\((218/3,54,17/3)\)\\
\(d_4\)&8&\((105/2,237/4,49/4)\)&\((149/2,54,23/4)\)\\
\(u_1\)&4&\((65,66,9)\)&\((80,54,6)\)\\
\(u_3\)&12&\((205/3,119/2,113/12)\)&\((74,54,23/4)\)\\
\bottomrule
\end{longtable}

\subsection{Secteur nul et statut des réalisations physiques}

\begin{theorem}[Bifurcation nulle]
\label{10j-thm-sector-nulo}
La carte photonique \(\mathcal N_\gamma\) et la carte gluonique \(\mathcal N_g\) bifurquent avant le caractère positif. Dans celles-ci
\begin{equation}
 \kappa_{\rm null}=\mathrm{N/A},
 \qquad m_\gamma=0,
 \qquad m_g=0,
 \label{10j-eq-nulidad-gauge}
\end{equation}
sans que l'extension, l'énergie ou la représentation de jauge soient triviales. La nullité ne s'obtient pas en imposant \(\kappa=0\) dans \(R_{\rm act}^\kappa\).
\end{theorem}

\begin{proof}
Le caractère positif prend ses valeurs dans \(\mathbb R_{>0}\), de sorte que \(R_{\rm act}^0=1\), et non zéro. Aucune puissance dans cette carte ne peut donc produire la section nulle. La réalisation photonique réside sur le bord nul électromagnétique et laisse deux polarisations transverses après la réduction de jauge ; la réalisation gluonique réside dans la représentation adjointe de couleur de dimension huit. L'inclusion des deux cartes dans la section zéro de \(\mathcal L_M\) s'effectue avant \(\chi_0\), ce qui prouve \eqref{10j-eq-nulidad-gauge} sans effacer leurs représentations.
\end{proof}

Les réalisations physiques ultérieures sont ordonnées selon leur domaine, leur représentation et leur modalité de publication ; leur nombre documentaire n'est pas le cardinal de l'atlas et ne définit pas la Loi générale des masses. Leur statut exact est :
\begin{longtable}{@{}p{.22\textwidth}p{.27\textwidth}p{.43\textwidth}@{}}
\caption{Statut probatoire des réalisations physiques ultérieures.}
\label{10j-tab-realizaciones-fisicas}\\
\toprule
classe(s)&salida interna&statut de la coordonnée de masse\\
\midrule
\endfirsthead
\toprule
classe(s)&salida interna&statut de la coordonnée de masse\\
\midrule
\endhead
électron&fibre \(\ell_0\), deux lecteurs coïncidents&scalaire bêta exact, prospectif et sans cible, \eqref{10j-eq-electron-beta}\\
muon, tau&opérateurs des fibres leptoniques et évaluations de signature conservées&opérateur exact ; les évaluations scalaires héritées sont conditionnées par la signature, et non des sélecteurs prospectifs\\
tres neutrinos&quatre fibres neutriniques, limites et opérateur de mélange&aucune masse de saveur n'est inventée ; la publication primaire est opératorielle\\
seis quarks&fibres \(u_j,d_j\), génération et orbite ternaire de couleur&opérateur exact ; tout scalaire de comparaison doit déclarer le schéma et l'échelle de renormalisation\\
photon&section nulle électromagnétique à deux modes transverses&nullité exacte \(m_\gamma=0\) dans sa carte\\
gluons&section nulle dans l'adjointe de couleur, avec huit générateurs&nullité exacte \(m_g=0\) dans sa carte ; l'écart collectif n'est pas une masse de pôle du gluon\\
\(W,Z\)&opérateurs de jauge massifs et signatures conditionnées&l'évaluation conditionnée est conservée ; la jointure voie--spectre physique n'est pas remplacée par la valeur externe\\
Higgs
&opérateur du secteur scalaire&scalaire physique conditionné par l'image directe de représentation\\
méson, baryon&grammaire de composition des registres et opérateur de liaison&ce ne sont pas des sommes nues de masses constituantes\\
Fibonacci
&anneau de fusion et représentation de tressage&codomaine topologique ; il n'y a pas de scalaire universel de masse\\
Ising/Majorana&anneau de fusion, autoconjugaison et terme compatible&la neutralité n'implique à elle seule ni Majorana ni masse\\
\(\mathbb Z_6\)&secteur topologique d'ordre six&opérateur dans son codomaine ; il n'est pas identifié à une hauteur de tour\\
virtuels&sections de persistance finie&classification par horizon terminal, non par masse sur couche\\
\bottomrule
\end{longtable}

\subsection{Complétude du spectre interne et de l'inventaire comparatif}

Le mot \emph{complet} s'applique ici à deux objets distincts et tous deux sont expressément déclarés. L'espace ambiant des deux spectres marginaux est
\begin{equation}
 \mathscr A_M:=
 \coprod_{M\in\mathcal M_{13}}
 \left(
  \operatorname{Spec}\widehat M_M^{\beta,D}
  \times
  \operatorname{Spec}\widehat M_M^{\beta,\Omega}
 \right)
 \ \coprod\ \{0_\gamma,0_g\},
 \label{10j-eq-espectro-interno-completo}
\end{equation}
avec conservation des multiplicités, projecteurs, voies, schéma et échelle. La sortie intrinsèque n'est toutefois pas tout ce produit : c'est le support de la mesure spectrale conjointe, définie dans \eqref{10j-eq-espectro-conjunto-completo}, qui conserve quels deux éléments proviennent d'une même voie. L'inventaire externe complet utilisé pour la reconnaissance ultérieure est
\begin{equation}
 \mathscr I_{\rm ext}
 =\mathscr I_m^{471}\coprod\mathscr I_\Gamma^{384},
 \qquad |\mathscr I_{\rm ext}|=471+384=855.
 \label{10j-eq-inventario-471-384}
\end{equation}
Chaque enregistrement de \(\mathscr I_m^{471}\) distingue masse stable, masse de pôle ou estimateur de résonance ; chaque enregistrement de \(\mathscr I_\Gamma^{384}\) distingue la largeur extraite d'un pôle ou d'un mode ouvert. La correspondance métrologique est la relation typée
\begin{equation}
 \mathcal C_{\rm met}
 \subset
 \mathscr A_M\times\mathscr I_{\rm ext},
 \label{10j-eq-correspondencia-metrologica-completa}
\end{equation}
qui ne s'ouvre qu'après le gel du descripteur, de la fibre, de l'opérateur, du projecteur, de la représentation, du schéma et de l'échelle.

\begin{theorem}[Double complétude sans sélection rétrospective]
\label{10j-thm-doble-completitud}
La Loi générale des masses publie une sortie opératorielle sur chacune des treize multisections et sur les (324) voies, en plus des deux sections nulles. L'inventaire de \eqref{10j-eq-inventario-471-384} individualise, sans omission documentaire, les (855) observables externes et conserve pour chacun son secteur HMT d'arrivée et la réalisation structurelle requise. Aucune grandeur de cet inventaire ne sélectionne une cellule, une voie, une valeur propre, une tour ou un coefficient de la loi.
\end{theorem}

\begin{proof}
Le \cref{10j-thm-dominio} épuise l'atlas et le \cref{10j-thm-ley-operatoria} construit les deux opérateurs positifs sur chaque fibre ; le \cref{10j-thm-sector-nulo} ajoute les deux cartes de masse nulle. Ainsi, le support conjoint de \eqref{10j-eq-espectro-conjunto-completo} couvre le domaine générateur intégral et est inclus dans l'espace ambiant de \eqref{10j-eq-espectro-interno-completo} sans lui être identifié. L'inventaire matériel énumère exactement (471) enregistrements de masse et (384) de largeur et attribue à chaque ligne une identité, une unité, une incertitude ou une limite, un schéma, une échelle, un secteur et une réalisation. Permuter ou remplacer les grandeurs externes agit uniquement sur le second facteur de \eqref{10j-eq-correspondencia-metrologica-completa} ; cela laisse invariants \(\mathcal C_{81}\), \(\mathcal R_{324}\), les lecteurs, les tours et les opérateurs. Cela prouve simultanément l'exhaustivité et l'antériorité causale.
\end{proof}

\subsection{CKM et PMNS comme transports distincts}

Dans le secteur quark, l'incidence N42 produit
\begin{equation}
 M_{\rm CKM}=
 \begin{pmatrix}
 2&-5&28\\
 0&7&-25/2\\
 1&-20&-2/3
 \end{pmatrix},
 \qquad \det M_{\rm CKM}=-\frac{3857}{6},
 \label{10j-eq-matriz-ckm}
\end{equation}
et transporte l'opérateur down dans la base up au moyen de
\begin{equation}
 B_d^{(u)}=V_{\rm CKM}B_dV_{\rm CKM}^*,
 \qquad
 R_{\rm act}^{B_d^{(u)}}
 =V_{\rm CKM}R_{\rm act}^{B_d}V_{\rm CKM}^*.
 \label{10j-eq-transporte-ckm}
\end{equation}
Le transport change la base, non les valeurs propres, et agit après que les lecteurs ont construit \(B_u,B_d\). La même matrice d'incidence agit sur les coordonnées angulaires déjà construites au moyen de
\begin{equation}
 \begin{pmatrix}\theta_{12}\\\theta_{23}\\\theta_{13}\end{pmatrix}
 =M_{\rm CKM}
 \begin{pmatrix}
  A\\ C^*/6\\ (180/\pi_{\rm HMT})\Delta_4
 \end{pmatrix},
 \qquad \delta_{\rm CP}=9A,
 \label{10j-eq-angulos-ckm}
\end{equation}
d'où
\begin{equation}
 (\theta_{12},\theta_{23},\theta_{13},\delta_{\rm CP})
 =(13.003313589509,\,2.398056157332,\,0.213977382244,
 \,65.676173123554)^\circ
 \label{10j-eq-salida-ckm}
\end{equation}
et
\begin{equation}
 J_{\rm HMT}=3.1189723326632974\times10^{-5},\qquad
 3857\delta_{\rm CP}
 =9(1528\theta_{12}+3380\theta_{23}+801\theta_{13}).
 \label{10j-eq-falsador-ckm}
\end{equation}
La dernière identité est un réfutateur linéaire exact du transport, non un ajustement de ses angles.

Pour le secteur leptonique, soient \(F_\ell,F_\nu:\mathbb C^3\to\mathcal H\) des isométries de même image \(\operatorname{im}P_3\), soit \(b_K\in\operatorname{im}P_3\) unitaire et soit \(Q_K=|b_K\rangle\langle b_K|\leq P_3\) son projecteur de rang un. Définissons
\begin{equation}
 \Phi_L=I+(e^{i\pi/6}-1)Q_K,
 \qquad U_{\rm HMT}:=F_\ell^*\Phi_LF_\nu.
 \label{10j-eq-unitaria-leptonica}
\end{equation}

\begin{proposition}[Distinction et complétude des transports CKM--PMNS]
\label{10j-prop-ckm-pmns}
La matrice \(U_{\rm HMT}\) de \eqref{10j-eq-unitaria-leptonica} est unitaire, mais la phase de rang un ne constitue pas la réalisation physique PMNS de rang plein. Le transport PMNS exact est défini par les repères spectraux complets, \(U_{\rm PMNS}=F_\ell^*F_\nu\). Le noyau \(G_0\) de registre fournit un certificat fini supplémentaire si et seulement si ses entrées sont présentées et si \(\det G_0\ne0\) est vérifié ; la condition n'est pas érigée ici en fait sans ce calcul. CKM et PMNS ont des domaines, des bases et des réfutateurs distincts et ne peuvent être identifiés.
\end{proposition}

\begin{proof}
Comme \(Q_K\leq P_3\) est un projecteur orthogonal, \([\Phi_L,P_3]=0\) et \(\Phi_L\) est unitaire. De plus, \(F_\ell F_\ell^*=F_\nu F_\nu^*=P_3\) ; par conséquent
\[
 U_{\rm HMT}^*U_{\rm HMT}
 =F_\nu^*\Phi_L^*P_3\Phi_LF_\nu=I_3.
\]
Cependant, une phase de rang un conserve un stabilisateur \(U(2)\). Le test de ses trente-six permutations ne reproduit pas le petit angle de mélange physique. L'échec de la phase minimale élimine un substitut dégénéré, mais n'ouvre ni ne rabaisse la réalisation PMNS, car les repères complets produisent directement \(F_\ell^*F_\nu\). Le facteur polaire d'un noyau de registre n'est un certificat alternatif que lorsqu'il satisfait son critère de rang. CKM, en revanche, transporte entre les bases up/down d'opérateurs quark déjà construits au moyen de \eqref{10j-eq-transporte-ckm}. Cette différence de domaines empêche de fusionner les deux constructions.
\end{proof}

\subsection{Fermeture spectrale conjointe et statut exact des transports CKM--PMNS}
\label{10j-sec-cierre-espectral-mezclas}

\paragraph{Spectre conjoint avec corrélation de voie.}
Pour chaque multisection \(M\), l'étape canonique antérieure à tout couplage non diagonal publie la paire
\begin{equation}
 \widehat M_{M,\mathrm{rut}}^{\beta,r}
 :=\sum_{\Gamma\in M}m_r(\Gamma)
 |e_\Gamma\rangle\langle e_\Gamma|,
 \quad
 m_r(\Gamma):=m_{\rm clk}\,
 \chi_{\rm full}^{\,r}\bigl(\sigma(\Gamma),\eta^r(\Gamma)\bigr)
 R_{\rm act}^{\kappa_r(\Gamma)},
 \quad r\in\{D,\Omega\}.
 \label{10j-eq-par-diagonal-ruta}
\end{equation}
Les deux opérateurs sont positifs et commutent. Il existe donc une unique PVM conjointe \(E_M^{D,\Omega}\) sur \(\mathbb R_+^2\) telle que
\begin{equation}
 \widehat M_{M,\mathrm{rut}}^{\beta,D}=\int x\,dE_M^{D,\Omega}(x,y),
 \qquad
 \widehat M_{M,\mathrm{rut}}^{\beta,\Omega}=\int y\,dE_M^{D,\Omega}(x,y),
 \label{10j-eq-pvm-conjunta}
\end{equation}
\begin{equation}
 E_M^{D,\Omega}(A)e_\Gamma
 =\mathbf 1_A\!\bigl(m_D(\Gamma),m_\Omega(\Gamma)\bigr)e_\Gamma.
 \label{10j-eq-pvm-accion-ruta}
\end{equation}
La mesure image
\begin{equation}
 \nu_M:=\frac1{|M|}\sum_{\Gamma\in M}
 \delta_{(m_D(\Gamma),m_\Omega(\Gamma))}
 \label{10j-eq-medida-imagen-conjunta}
\end{equation}
conserve la corrélation entre les deux lecteurs et toutes leurs multiplicités. En conséquence, le produit cartésien écrit dans \eqref{10j-eq-espectro-interno-completo} est l'espace ambiant ; l'objet spectral intrinsèque est
\begin{equation}
 \begin{aligned}
 \mathscr A_M^{\rm joint}
 &:=
 \coprod_{M\in\mathcal M_{13}}\operatorname{supp}E_M^{D,\Omega}
 \ \coprod\ \{0_\gamma,0_g\}\\
 &\hookrightarrow
 \coprod_{M\in\mathcal M_{13}}
 \bigl(\operatorname{Spec}\widehat M_M^{\beta,D}
 \times\operatorname{Spec}\widehat M_M^{\beta,\Omega}\bigr)
 \ \coprod\ \{0_\gamma,0_g\}.
 \end{aligned}
 \label{10j-eq-espectro-conjunto-completo}
\end{equation}
L'inclusion n'est pas érigée en égalité sans prouver que toutes les combinaisons marginales appartiennent à une même voie. Pour un canal préparé \(P_{X,a}\), la compression \(P_{X,a}E_M^{D,\Omega}(\cdot)P_{X,a}\) est la POVM induite. Il existe une ligne de masse \(m_X\) seulement si
\begin{equation}
 E_{\mathcal M_X}(\{m_X\})P_{X,a}=P_{X,a},
 \qquad
 (\mathcal M_X-m_XI)P_{X,a}=0.
 \label{10j-eq-criterio-linea-espectral}
\end{equation}
Une compression comportant plusieurs points de support publie un multiplet ou une distribution, non une moyenne. Si un couplage ultérieur cesse d'être diagonal, la paire ordonnée est conservée et une PVM conjointe n'est affirmée qu'après vérification du commutateur nul.

\paragraph{Transporte CKM.}
Avec \(s_{ij}=\sin\theta_{ij}\), \(c_{ij}=\cos\theta_{ij}\) et \(\delta=\delta_{\rm CP}\), la sortie de \eqref{10j-eq-angulos-ckm} détermine
\begin{equation}
 V_{\rm CKM}=\begin{pmatrix}
 c_{12}c_{13}&s_{12}c_{13}&s_{13}e^{-i\delta}\\
 -s_{12}c_{23}-c_{12}s_{23}s_{13}e^{i\delta}&
 c_{12}c_{23}-s_{12}s_{23}s_{13}e^{i\delta}&s_{23}c_{13}\\
 s_{12}s_{23}-c_{12}c_{23}s_{13}e^{i\delta}&
 -c_{12}s_{23}-s_{12}c_{23}s_{13}e^{i\delta}&c_{23}c_{13}
 \end{pmatrix}.
 \label{10j-eq-vckm-explicita}
\end{equation}
Elle est unitaire, réalise le changement de repère de \eqref{10j-eq-transporte-ckm} et ne crée pas les valeurs propres. Si \(\kappa_i^u\) et \(\kappa_a^d\) sont les spectres des opérateurs hermitiens de type haut et bas, alors
\begin{equation}
 \left|\det[B_u,V_{\rm CKM}B_dV_{\rm CKM}^*]\right|
 =2|J_{\rm HMT}|
 \prod_{i<j}|\kappa_i^u-\kappa_j^u|
 \prod_{a<b}|\kappa_a^d-\kappa_b^d|.
 \label{10j-eq-determinante-conmutador-ckm}
\end{equation}
C'est le réfutateur interne spectre--mélange. L'inférence statistique face aux données externes reste dans la couche de reconnaissance et exige leur covariance.

\paragraph{Transport PMNS et critère de rang plein.}
La phase minimale de \eqref{10j-eq-unitaria-leptonica} satisfait \(Q_K\leq P_3\), donc \([\Phi_L,P_3]=0\) ; son unitarité est exacte, bien qu'elle demeure un contrôle négatif de rang un. Soient \(M_\ell\) et \(M_\nu\) les opérateurs chargé et neutre déjà construits, et soient \(F_\ell,F_\nu\) des repères spectraux orthonormaux complets, avec \(F_\ell F_\ell^*=F_\nu F_\nu^*=P_3\). Le transport spectral est
\begin{equation}
 U_{\rm PMNS}:=F_\ell^*F_\nu,
 \qquad U_{\rm PMNS}^*U_{\rm PMNS}=I_3.
 \label{10j-eq-pmns-marcos-espectrales}
\end{equation}
Si les deux spectres sont simples, il est déterminé à des réajustements diagonaux de phase et à l'ordre spectral déclaré près. S'il existe des dégénérescences, la sortie exacte est la double classe
\begin{equation}
 \prod_\lambda U(m_{\ell,\lambda})\backslash U(3)/
 \prod_\mu U(m_{\nu,\mu}),
 \label{10j-eq-pmns-clase-doble}
\end{equation}
non une matrice numérique choisie dans les sous-espaces dégénérés.

Le noyau de registre conserve la retenue, le feuillet, l'orientation, la lacune et l'holonomie. En omettant les termes de norme nulle, on peut écrire
\begin{equation}
 \begin{aligned}
 K_{\rm reg}(c,d)&:=
 \sum_{r\in\{\mathrm{car},\mathrm{hoj},\mathrm{ori},\mathrm{vac},\mathrm{hol}\}}
 \frac{\langle\Xi_r(c),\Xi_r(d)\rangle}{\sum_x\|\Xi_r(x)\|^2},\\
 (G_0)_{ab}&:=\frac1{\sqrt{|L_a||N_b|}}
 \sum_{c\in L_a}\sum_{d\in N_b}K_{\rm reg}(c,d).
 \end{aligned}
 \label{10j-eq-pmns-kernel-completo}
\end{equation}
La condition
\begin{equation}
 \det G_0\ne0
 \quad\Longleftrightarrow\quad
 \text{le registre déclaré sépare trois directions},
 \qquad
 U_{G_0}:=G_0(G_0^*G_0)^{-1/2},
 \label{10j-eq-pmns-criterio-rango}
\end{equation}
est un critère fini exact. Il n'est pas affirmé comme fait certifié tant que les entrées de \(G_0\) et le calcul de son déterminant ne sont pas présentés. Si \(\det G_0=0\), il faut ajouter un observable interne indépendant et répéter le calcul ; il est interdit de forcer l'inversibilité au moyen d'angles PMNS externes. La PVM conjointe, sa mesure image, le transport CKM et le théorème des repères PMNS sont des résultats internes exacts avec la structure de départ déclarée ; le critère de \(G_0\) reste distinct de ces résultats.

\subsection{Non-gouvernance des tableaux historiques}

\begin{proposition}[Indépendance prospective]
\label{10j-prop-tablas-historicas}
Les anciennes fenêtres de dix signatures et douze masses sont des témoins historiques de comparaison. Elles ne définissent ni \(\mathcal C_{81}\), ni \(\mathcal F_{56}\), ni \(\mathcal M_{13}\), ni \(\mathcal R_{324}\), ni les lecteurs, ni le caractère, ni les tours. Permuter, supprimer ou remplacer leurs colonnes métrologiques laisse inchangées toutes les constructions de \eqref{10j-eq-cadena-masas} antérieures à la comparaison.
\end{proposition}

\begin{proof}
Les données d'entrée de la construction sont l'atlas, les quatre orientations, la transition \(U_t\), le registre \(81\times108\), la signature et les lecteurs. Aucune des définitions \eqref{10j-eq-rutas-324}, \eqref{10j-eq-firma-z5}, \eqref{10j-eq-kd}, \eqref{10j-eq-komega} ou \eqref{10j-eq-semigrupo} ne contient de masse externe. Par conséquent, une permutation du catalogue ultérieur commute avec la projection qui oublie la comparaison et agit comme l'identité sur le constructeur interne. Un tableau abrégé ne peut donc être ni le domaine ni le générateur de la loi.
\end{proof}

\subsection{Théorème de fermeture typée et réfutateurs}

\begin{theorem}[Loi générale des masses HMT--MD]
\label{10j-thm-ley-general}
La chaîne \eqref{10j-eq-cadena-masas} définit sur le domaine \(81/56/13/324\), avec un secteur nul séparé, une loi de masse qui satisfait les propriétés suivantes :
\begin{enumerate}
 \item construit la voie, le registre, la signature, le caractère et la tour sans grandeurs externes ;
 \item fixe l'échelle absolue au moyen du conoyau d'ordre seize et de la décennie \(-34\), séparée de tout correcteur relatif ;
 \item publie deux lecteurs non équivalents avec les recensements \(14/9/40/18\) et l'unique coïncidence \((80,54,6)\) ;
 \item ferme l'électron par la réduction de rang un \eqref{10j-eq-electron-beta} ;
 \item étend exactement la loi sous la forme des paires d'opérateurs \((B_M^D,B_M^\Omega)\) sur les treize multisections ;
 \item fait bifurquer le photon et les gluons vers des cartes nulles avant le caractère ;
 \item publie sur chaque fibre non centrale l'opérateur positif, sa mesure spectrale complète et ses multiplicités ; toute coordonnée scalaire physique est l'image de l'image directe et de la scalarisation déclarées pour cette représentation ;
 \item incorpore, comme couche comparative ultérieure, l'inventaire exhaustif des (471) masses et (384) largeurs sans permettre à leurs valeurs de rétroagir sur l'atlas ;
 \item maintient CKM et PMNS distincts, construit la sortie angulaire CKM avec son réfutateur linéaire et le transport PMNS de rang plein, et n'ouvre PDG/CODATA qu'après le gel de la sortie interne.
\end{enumerate}
Ce n'est pas une loi d'une seule paire angulaire : \(A,C^*\) paramètrent les canaux globaux, tandis que la séparation des espèces réside dans la voie, le registre, les deux lecteurs, les tours, la représentation et le couplage.
\end{theorem}

\begin{proof}
Les points 1 et 2 sont les \cref{10j-lem-caracter,10j-thm-escala-menos34,10j-prop-escala-corrector} ; le point 3 est le \cref{10j-thm-censo-lectores} ; le point 4 est le \cref{10j-thm-electron} ; les points 5 et 6 sont les \cref{10j-thm-ley-operatoria,10j-thm-sector-nulo} ; le statut du point 7 est typé dans \cref{10j-tab-realizaciones-fisicas} ; le point 8 est le \cref{10j-thm-doble-completitud} ; et le point 9 est le \cref{10j-prop-ckm-pmns}. La \cref{10j-prop-tablas-historicas} prouve que la comparaison externe ne peut rétroagir sur cette chaîne. Ensemble, les constructions composent la loi affirmée sans identifier l'échelle, le caractère, le lecteur, l'opérateur et la scalarisation.
\end{proof}

\begin{falsifier}[Contrôles négatifs de la loi]
\label{10j-falsador-ley-masas}
La construction est réfutée ou mal typée si l'un des faits suivants se produit :
\begin{enumerate}
 \item une masse, CODATA, PDG ou un tableau 10/12 sélectionne une cellule, une voie, une signature, un coefficient ou une valeur propre ;
 \item le saut APP \(4\) est identifié au défaut, au déploiement ou au retour \(54\) ;
 \item on insère \(-34\) dans \(\kappa\) ou l'on multiplie de nouveau une coordonnée MeV par \(10^{-34}\) ;
 \item les recensements des lecteurs cessent d'être \(14,9,40,18\), ou une coïncidence autre que \((80,54,6)\) apparaît ;
 \item la valeur basale \(e_0\) est présentée comme sortie électronique finale ;
 \item le semi-groupe \(\langle90,120\rangle\) est réduit à une liste finie ou utilisé pour sélectionner rétrospectivement une espèce ;
 \item une fibre de rang supérieur à un reçoit un scalaire sans image directe physique ni scalarisation prospectives ;
 \item la nullité du photon ou du gluon est déduite de \(R_{\rm act}^{0}=1\), ou l'écart collectif est confondu avec une masse de pôle gluonique ;
 \item CKM et PMNS sont identifiés ou la phase leptonique de rang un est érigée en mélange physique complet ;
 \item une ligne de l'inventaire (471+384) manque, une largeur est confondue avec une masse ou une valeur externe modifie le spectre interne ;
 \item la chaîne généalogique est remplacée par une unique paire angulaire.
\end{enumerate}
\end{falsifier}

L'architecture de la loi est une \textsc{architecture auctoriale préexistante} ; la récupération conjointe de l'atlas, des lecteurs et de l'électron est un \textsc{résultat récupéré} ; la présentation opératorielle sur les treize fibres est une formalisation réunie ; et les recensements exécutables sont des certificats. Aucune de ces catégories n'altère la précédence causale d'APP--TRIT--TPK.
